%%%%%%%%%%%%%%%%%%%%%%%%%%%%%%%%%%%%%%%%%%%%%%%%%%%%%%%%%%%%%%%%%%
%% Toto je inkludovaný LaTeXový soubor `ca_mirr.tex'.           %%
%% ------------------------------------------------------------ %%
%% Verze 2022-05-11                                             %%
%% Vyvíjí & udržuje <stanislav.hledik@physics.slu.cz>           %%
%%%%%%%%%%%%%%%%%%%%%%%%%%%%%%%%%%%%%%%%%%%%%%%%%%%%%%%%%%%%%%%%%%

\chapter{Umístění sazebního obrazce}\label{mirror}%%%%%%%%%%%%%%%%%%%%%%%%%

Implicitně je sazební obrazec horizontálně centrován a jeho vertikální pozice
je optimalizována pro tisk na papír
formátu~A4~\cite{SSMGT:LIVE::FormatyPapiru}. Může se však stát, že uživatel
bude z~nějakého důvodu chtít posunout sazební obrazec například směrem dovnitř
(ke hřbetu), například proto, že pro svázání práce použije lepenou vazbu, nebo
směrem ven (od hřbetu) v~případě obvyklé vazby nabízené většinou
knihařství. Také může být žádoucí posunout sazební obrazec ve vertikálním
směru, například kvůli ořezu listů.  K~tomu účelu slouží příkazy
\pgm{\textbackslash spineoffset} a \pgm{\textbackslash verticalshift}.

Příkaz \pgm{\textbackslash spineoffset} posune sazební obrazec o~délku
specifikovanou argumentem směrem ke hřbetu (tj. sudé stránky, které jsou při
rozevření nalevo, budou posunuty o~hodnotu argumentu doprava a liché, které
jsou při rozevření napravo, budou posunuty o~hodnotu argumentu doleva). To
platí pro kladný argument, např.~\pgm{5mm}. Zadáme\discretionary{-}{-}{-}li
záporný argument, např.~\pgm{-24pt}, bude sazební obrazec posunut směrem ven,
od hřbetu.

Podobně funguje příkaz \pgm{\textbackslash verticalshift}. Kladná hodnota
argumentu posouvá sazební obrazec nahoru, záporná dolů. V~obou příkazech lze
použít libovolné jednotky akceptované \TeX em~\cite{Ryb:2005:LaTeX:3,%
  Kop-Dal:2004:LaTeXPodrPruv:,%
  Goo-Mit-Sam:1994:Companion:}. Samotné rozměry sazebního obrazce zůstanou
nedotčeny. Příkazy pro adjustaci polohy sazebního obrazce \pgm{\textbackslash
  spineoffset} a \pgm{\textbackslash verticalshift} lze výhodně umístit do
konfiguračního souboru \pgm{\ipotname .cfg}, např.

{\small
\begin{verbatim}
%% Page positioning
\spineoffset{-5mm}
\verticalshift{2mm}
\end{verbatim}
\par}

\noindent
Popsané chování je nezávislé na aktivaci volby \opt{narrowmargin} (viz
Dodatek~\ref{options} na stránce~\pageref{opnarrowmargin}). Je pouze třeba
uvážit, že při aktivaci této volby jsou okraje užší a rozsah použitelných
posuvů odpovídajícím způsobem menší.

Při umisťování sazebního obrazce je vhodné řídit se pravidly estetické
sazby, např. podle~\cite{Pop-Fle-Pop:1984:SazbaI:}.

\endinput

%%
%% End of file `ca_mirr.tex'.
